% GENERATED FILE. Edit canonical lessons, then run npm run generate:books.
% canonical-source-hash: fnv1a64:1d5690ed4a09f9cf
% canonical-lessons: IT-C35-sapere, IT-C35-non-so, IT-C35-oggi-ieri, IT-C35-molto-poco, IT-C35-abbastanza

\chapter{Knowing, and How Much}
\label{ch:it-sapere}
\hlchaptermodality{35}

\begin{chapteropening}
\textbf{By the end of this chapter:} \emph{I can say that I know a fact and that I know a person, say that I do not know, place today and yesterday beside tomorrow, and put something on the scale from molto through abbastanza to un po'.}

The chapter's last lesson completes the quantity scale and gives Come stai? — a question from Chapter 2 — its middle answer: abbastanza bene.
\end{chapteropening}

\section[sapere]{sapere — to know, and to know how}

\label{lesson:IT-C35-sapere}



\begin{quote}
A new chapter. Five retrievals, then a verb this book has gone without.

\begin{itemize}
\raggedright
  \item \emph{Recall:} say \emph{per me, a te} — \textbf{R1}, one lesson back, before it decays
  \item \emph{Recall:} say \emph{mio, tuo, suo} — \textbf{R2}, the first real retrieval, five to fifteen lessons back
  \item \emph{Recall:} say \emph{che} — \textbf{R3}, durable, twenty to sixty lessons back
  \item \emph{Recall:} say \emph{essere} — \textbf{R4}, recognition at distance, eighty lessons back or more
\end{itemize}
\end{quote}

\subsection*{You'll want to know: the verb, and its irregular first person}
\begin{quote}
\textbf{sapere} — \textbf{to know} \emph{(a fact)}, and \textbf{to know how} \emph{(to do something)}
\end{quote}

It is irregular, and the three singular forms are the ones to have:

\noindent
\begin{tabularx}{\linewidth}{@{}>{\raggedright\arraybackslash}X>{\raggedright\arraybackslash}X@{}}
\toprule
\textbf{} & \textbf{} \\
\midrule
\textbf{so} & I know \\
\textbf{sai} & you know \\
\textbf{sa} & he / she knows \\
\bottomrule
\end{tabularx}

\textbf{So} is the shortest verb form in the language. The plurals are \emph{sappiamo, sapete, sanno}.

\begin{grammarlens}[title={Grammar lens: sapere plus an infinitive}]
Put an infinitive straight after it and \emph{sapere} stops meaning \emph{know} and starts meaning \textbf{can, know how to}:

\begin{quote}
\textbf{So parlare italiano.} — I can speak Italian. \emph{(I have learnt how)} \textbf{Sai leggere?} — Can you read?
\end{quote}

No \emph{per} and no \emph{to}: just \textbf{sapere} and the plain infinitive. \emph{So parlare} says what \emph{parlo} does not — that you learnt it.

\textbf{Sapere} is Latin \textbf{sapere}, \textquotedblleft{}to taste.\textquotedblright{} A person of good taste became a person of good judgement, and judgement became knowledge — English \textbf{savour}, \textbf{insipid}, and \textbf{homo sapiens}, the tasting animal.

\textbf{And now the line English does not draw.} Italian has a second verb for knowing, and you already have it: the \textbf{incontrare} lesson gave you \textbf{conoscere} as \emph{incontrare}'s partner, traced it to Latin \emph{cognōscere}, and put it in the past. It was never set against \emph{sapere}, because \emph{sapere} did not exist in this book. Here is the line:

\begin{quote}
\textbf{sapere} — a \textbf{fact}, or how to do something: \emph{so parlare italiano} \textbf{conoscere} — a \textbf{person} you have met: \emph{conosco Marco}
\end{quote}

\emph{So Marco} is wrong. Spanish draws the same line with \emph{saber} and \emph{conocer}, French with \emph{savoir} and \emph{connaître}. English is the odd one out.
\end{grammarlens}

\subsection*{Your turn}
Say these aloud:

\begin{itemize}
\raggedright
  \item \emph{so, sai, sa}
  \item \emph{So parlare italiano.}
  \item \emph{Sai leggere?}
  \item \emph{Conosco Marco.}
\end{itemize}

\begin{tcolorbox}[breakable,colback=teal!4,colframe=teal!35!black,title={Before you move on}]
What is the \emph{io} form of \emph{sapere}? (\emph{\textbf{So}}.) What does \emph{sapere} plus an infinitive mean? (\textbf{To know how} — \emph{so parlare}.) And which verb takes a person instead? (\emph{\textbf{Conoscere}}, which the \emph{incontrare} lesson already gave you.)
\end{tcolorbox}
\section[non so]{non so — I don't know}

\label{lesson:IT-C35-non-so}



\begin{quote}
One lesson back, and the verb that made this possible.

\begin{itemize}
\raggedright
  \item \emph{Recall:} say \emph{sapere} — \textbf{R1}, one lesson back, before it decays
  \item \emph{Recall:} say \emph{nostro, vostro} — \textbf{R2}, the first real retrieval, five to fifteen lessons back
  \item \emph{Recall:} say \emph{penso che} — \textbf{R3}, durable, twenty to sixty lessons back
  \item \emph{Recall:} say \emph{essere $\to$ stato} — \textbf{R4}, recognition at distance, eighty lessons back or more
\end{itemize}
\end{quote}

\subsection*{You'll want to know: three things at once}
\begin{quote}
\textbf{Non so.} — I don't know.
\end{quote}

\textbf{This lesson introduces no new word.} \emph{Non} is the negation chapter's, \emph{so} is one lesson back, and the rule that \emph{non} goes immediately before the verb does the rest. It is worth its own lesson because those two syllables do three different jobs.

\textbf{One — plain ignorance.} Somebody asks you something you do not know:

\begin{quote}
— \textbf{Chi è?} — \textbf{Non so.}
\end{quote}

\textbf{Two — hedging.} Put it in front of a question word and it becomes a way of declining to commit:

\begin{quote}
\textbf{Non so quando.} — I don't know when. \textbf{Non so perché.} — I don't know why. \textbf{Non so chi è.} — I don't know who it is.
\end{quote}

Every one of those question words is Chapters 31 and 32's, and \emph{non so} is what lets you use them about yourself rather than only to ask.

\textbf{Three — scepticism.} A bare, slow \textbf{Non so …} with the sentence left hanging is how an Italian expresses doubt about what has just been said. It is the politest disagreement in the language, and it is a great deal gentler than \emph{non sono d'accordo}.

The pair to hold together is the pair a beginner needs hourly:

\begin{quote}
\textbf{Non capisco.} — I didn't follow you. \textbf{Non so.} — I followed you, and I have no answer.
\end{quote}

\subsection*{Your turn}
Say these aloud:

\begin{itemize}
\raggedright
  \item \emph{Non so.}
  \item \emph{Non so quando.}
  \item \emph{Non so perché.}
  \item \emph{Non so chi è.}
\end{itemize}

\begin{tcolorbox}[breakable,colback=teal!4,colframe=teal!35!black,title={Before you move on}]
Say \textquotedblleft{}I don't know.\textquotedblright{} (\emph{\textbf{Non so}}.) Say \textquotedblleft{}I don't know why.\textquotedblright{} (\emph{\textbf{Non so perché}}.) What is the difference between \emph{non so} and \emph{non capisco}? (\emph{\textbf{Non capisco}} \textbf{means I did not follow;} \emph{\textbf{non so}} \textbf{means I followed and have no answer}.)
\end{tcolorbox}
\section[oggi, ieri]{oggi, ieri — today and yesterday}

\label{lesson:IT-C35-oggi-ieri}



\begin{quote}
Two lessons in, and the throw back reaches the going verb.

\begin{itemize}
\raggedright
  \item \emph{Recall:} say \emph{non so} — \textbf{R1}, one lesson back, before it decays
  \item \emph{Recall:} say \emph{il mio nome} — \textbf{R2}, the first real retrieval, five to fifteen lessons back
  \item \emph{Recall:} say \emph{chi} — \textbf{R3}, durable, twenty to sixty lessons back
  \item \emph{Recall:} say \emph{andare} — \textbf{R4}, recognition at distance, eighty lessons back or more
\end{itemize}
\end{quote}

\subsection*{You'll want to know: the set, completed}
\begin{quote}
\textbf{oggi} — today · \textbf{ieri} — yesterday · \textbf{domani} — tomorrow
\end{quote}

You have had \textbf{domani} for a long time and only as the second half of a farewell, \textbf{a domani} — \textquotedblleft{}until tomorrow.\textquotedblright{} It was two thirds of a set with the other two thirds missing, so a learner asked when something happened had one word and it pointed forward.

Now the set works, and it works with the two pasts this book already taught:

\begin{quote}
\textbf{Oggi lavoro a Roma.} — Today I'm working in Rome. \textbf{Ieri ho lavorato a Milano.} — Yesterday I worked in Milan. \textbf{Domani prendo un caffè.}
\end{quote}

\emph{Oggi} and \emph{ieri} go where any adverb of time goes: at the front for emphasis, or after the verb. Both are invariable — they never change their ending.

\begin{cousinweb}
\textbf{Oggi} is Latin \textbf{hodiē}, and \emph{hodiē} is itself two words squashed together: \textbf{hōc diē}, \textquotedblleft{}on this day.\textquotedblright{} The \textbf{diē} in it is \emph{diēs}, \textquotedblleft{}day\textquotedblright{} — the very word behind \textbf{giorno}, which the first chapter traced through \emph{diurnum}. So \emph{oggi} and \emph{giorno} share a root, and \emph{oggi} is literally \emph{this-day}.

Every Romance language did the same and hid the seam differently: Spanish \emph{hoy}, French \emph{aujourd'hui} (which is \emph{au jour d'hui}, \textquotedblleft{}on the day of today\textquotedblright{}, with \emph{hui} being \emph{hodiē} worn to two letters and then reinforced all over again), Portuguese \emph{hoje}.

\textbf{Ieri} is Latin \textbf{herī}, and it is the same word as the first half of English \textbf{yesterday} — Old English \emph{ġeostran}, from the same ancient root. The word for yesterday has meant yesterday, with barely a change, for as long as anything can be traced.
\end{cousinweb}

\subsection*{Your turn}
Say these aloud:

\begin{itemize}
\raggedright
  \item \emph{oggi, ieri, domani}
  \item \emph{Oggi lavoro a Roma.}
  \item \emph{Ieri ho lavorato a Milano.}
\end{itemize}

\begin{tcolorbox}[breakable,colback=teal!4,colframe=teal!35!black,title={Before you move on}]
Name the three. (\emph{\textbf{Oggi, ieri, domani}}.) Which of them did this book already have, and where? (\emph{\textbf{Domani}}, \textbf{inside the farewell} \emph{\textbf{a domani}}.) What is \emph{oggi} made of? (\emph{\textbf{Hōc diē}}, \textquotedblleft{}on this day\textquotedblright{} — the same \emph{diēs} as \emph{giorno}.)
\end{tcolorbox}
\section[molto, poco]{molto, poco — a lot and a little}

\label{lesson:IT-C35-molto-poco}



\begin{quote}
Four in, and this one unfreezes two things you have been carrying.

\begin{itemize}
\raggedright
  \item \emph{Recall:} say \emph{oggi, ieri} — \textbf{R1}, one lesson back, before it decays
  \item \emph{Recall:} say \emph{signore, signora} — \textbf{R2}, the first real retrieval, five to fifteen lessons back
  \item \emph{Recall:} say \emph{che cosa} — \textbf{R3}, durable, twenty to sixty lessons back
  \item \emph{Recall:} say \emph{sono andato} — \textbf{R4}, recognition at distance, eighty lessons back or more
\end{itemize}
\end{quote}

\subsection*{You'll want to know: two jobs each}
\begin{quote}
\textbf{molto} — much, many, a lot, very · \textbf{poco} — little, few, not much
\end{quote}

Each of them does two jobs, and the jobs behave differently.

\textbf{In front of a describing word they mean \emph{very} and \emph{not very}, and they do not change:}

\begin{quote}
\textbf{Il caffè è molto buono.} — The coffee is very good. \textbf{Il tè è poco buono.} — The tea isn't very good.
\end{quote}

\textbf{In front of a noun they mean \emph{a lot of} and \emph{not much}, and they agree with it like any describing word:}

\begin{quote}
\textbf{molto caffè} · \textbf{molta acqua} · \textbf{molti amici} · \textbf{molte persone} \textbf{poco zucchero} · \textbf{poche persone}
\end{quote}

Before an adjective, frozen; before a noun, agreeing.

\emph{Molto} has never had a lesson. Where an Italian uses it most is the courtesy \textbf{molto piacere}, \textquotedblleft{}much pleasure\textquotedblright{} — the fuller form of the \emph{piacere} you say while shaking hands, and the first job above, doing its work inside a phrase nobody takes apart.

\begin{cousinweb}
The \emph{parlo un po'} dialogue said \textbf{Sì, parlo un po'} — \textquotedblleft{}yes, I speak a little\textquotedblright{} — and that \emph{un po'} was never glossed. Here is what it is.

\textbf{Po'} is \textbf{poco} with its last syllable cut off and an apostrophe left to mark the cut. Not an accent — an \textbf{apostrophe}, and it is compulsory, because \emph{un po} without it is not a word.

\begin{quote}
\textbf{un po'} — a little, a bit \textbf{un po' di zucchero} — a little sugar \textbf{Parlo un po' d'italiano.}
\end{quote}

\emph{Un po'} is softer than \emph{poco}. \textbf{Poco} is faintly negative — \emph{poco zucchero} is less than you wanted — while \textbf{un po'} is neutral and friendly, which is why that speaker chose it about their own Italian.

\textbf{Molto} is Latin \textbf{multus} (English \emph{multiple}); \textbf{poco} is Latin \textbf{paucus} (\emph{paucity}).
\end{cousinweb}

\subsection*{Your turn}
Say these aloud:

\begin{itemize}
\raggedright
  \item \emph{molto, poco}
  \item \emph{Il caffè è molto buono.}
  \item \emph{molti amici, molte persone}
  \item \emph{un po' di zucchero}
\end{itemize}

\begin{tcolorbox}[breakable,colback=teal!4,colframe=teal!35!black,title={Before you move on}]
When does \emph{molto} change its ending? (\textbf{Before a noun} — \emph{molti amici}. Before an adjective it is frozen.) What is \emph{un po'}? (\emph{\textbf{Poco}} \textbf{with its tail cut off} — and the mark is an apostrophe, not an accent.) Where had you already heard it? (\textbf{The} \emph{\textbf{parlo un po'}} \textbf{dialogue}.)
\end{tcolorbox}
\section[abbastanza]{abbastanza — enough, and quite}

\label{lesson:IT-C35-abbastanza}



\begin{quote}
The last lesson of the tranche, and it completes a set of three.

\begin{itemize}
\raggedright
  \item \emph{Recall:} say \emph{molto, poco} — \textbf{R1}, one lesson back, before it decays
  \item \emph{Recall:} say \emph{per me, a te} — \textbf{R2}, the first real retrieval, five to fifteen lessons back
  \item \emph{Recall:} say \emph{Che buono!} — \textbf{R3}, durable, twenty to sixty lessons back
  \item \emph{Recall:} say \emph{la testa} — \textbf{R4}, recognition at distance, eighty lessons back or more
\end{itemize}
\end{quote}

\subsection*{You'll want to know: the middle of the scale}
\begin{quote}
\textbf{abbastanza} — \textbf{enough}, and also \textbf{quite, fairly}
\end{quote}

It never changes its ending, whatever it stands in front of. That makes it the easiest of the three and the one to reach for when you are unsure:

\begin{quote}
\textbf{abbastanza caffè} — enough coffee \textbf{abbastanza persone} — enough people \textbf{Il caffè è abbastanza buono.} — The coffee is quite good.
\end{quote}

The three words now make a scale, and it is the scale an A1 conversation runs on:

\noindent
\begin{tabularx}{\linewidth}{@{}>{\raggedright\arraybackslash}X>{\raggedright\arraybackslash}X@{}}
\toprule
\textbf{} & \textbf{} \\
\midrule
\textbf{molto} & a lot / very \\
\textbf{abbastanza} & enough / quite \\
\textbf{poco} · \textbf{un po'} & not much / a little \\
\bottomrule
\end{tabularx}

And it answers a question you have been able to ask since the wellbeing chapter without being able to answer it in the middle: \textbf{Come stai?} — \textbf{Abbastanza bene.}

\begin{cousinweb}
\textbf{Abbastanza} is \textbf{a} + \textbf{bastanza}, built on the verb \textbf{bastare}, \textquotedblleft{}to be enough, to suffice.\textquotedblright{}

You know its third person singular already, even if you have never seen it in Italian, because English borrowed it whole: \textbf{basta!} — \textquotedblleft{}that's enough!\textquotedblright{} It is one of a handful of Italian words that travelled into other languages unchanged, alongside \emph{ciao} and \emph{piano}.

So \emph{abbastanza} is, literally, \textquotedblleft{}\textbf{to sufficiency}.\textquotedblright{} And \textbf{basta} is the same verb with nothing added: \emph{basta così} — that's enough like that — which uses \emph{così} one more time.
\end{cousinweb}

\subsection*{Your turn}
Say these aloud:

\begin{itemize}
\raggedright
  \item \emph{abbastanza}
  \item \emph{abbastanza caffè}
  \item \emph{Il caffè è abbastanza buono.}
  \item \emph{Abbastanza bene.}
\end{itemize}

\begin{tcolorbox}[breakable,colback=teal!4,colframe=teal!35!black,title={Before you move on}]
How do you say \textquotedblleft{}enough\textquotedblright{}? (\emph{\textbf{Abbastanza}}.) Does it change its ending? (\textbf{No} — never.) Which verb is it built on? (\emph{\textbf{Bastare}} — the \emph{basta!} English borrowed.) Answer \emph{Come stai?} from the middle of the scale. (\emph{\textbf{Abbastanza bene}}.)
\end{tcolorbox}
